\scanpage{13}
\deleted{2. Soit $X$ un espace, $Y$ un sous-espace fermé.}
\unsure{\deleted{(6)}|(6)} \deleted{Soit $T$}

\uline{NB} Les $T_\lambda$ $(\lambda\in[0,1])$ forment
une famille croissante de parties fermées de
$T$, les $\mathring T_\lambda\overset{\mathrm{def}}=
\lambda\mathring T\deleted{\unsure{_\lambda}}$ $(\lambda\in]0,1])$
\margin{$\Vert$}
une famille croissante de parties ouvertes. Si $\lambda\in]0,1[$, alors
$\mathring T_\lambda$ est bel et bien l'intérieur \inserted{topologique}
de $T_\lambda$ dans $T$,
mais bien sûr $\mathring T_1=\mathring T$ n'est pas l'intérieur
\inserted{top.} de $T_1$ dans
l'espace topologique $T_1=T$, i.e. on n'a pas
$\mathring T\unsure{\deleted{_1}|_1}=T$ sauf
si $\dot T=\phi$ i.e. $Y=T$.

Posons, pour $\lambda\in]0,1[$
\[
(\unsure{6})\qquad
\left\{\begin{aligned}
T\setminus\mathring T_\lambda
&\overset{\mathrm{def}}=\deleted{\unsure{E}}
\mathrm{Ep}^{(T)}_{1-\lambda=\varepsilon}
\quad\text{``épaississement d'indice $\varepsilon$''}\\
&\hspace{4em}\text{du bord $\dot T$ de $T$}\\
\dot{\mathrm{Ep}}_\varepsilon^{(T)}
&\overset{\mathrm{déf}}=\dot T\unsure{\amalg|\cup}(1-\varepsilon).\dot T
\end{aligned}\right.
\]
\uline{NB} La frontière de $E_\varepsilon$ dans $T$ est
$(1-\varepsilon)\dot T$, et $E_\varepsilon$
est un voisinage de $\dot T$. On a
\[
(7)\qquad\mathrm{Ep}_\varepsilon^{(T)}
\simeq[\underbrace{1-\varepsilon}_{\lambda},1]\times\dot T
\]
\margin{\uline{NB} Si $Y,\inserted{\dot T}$ compacts, \deleted{\unsure{les}}
alors $T$ compact, et les $T_\varepsilon$
$(\varepsilon\in]0,1])$ forment
un syst. fond. de voisinages de $Y$ dans $T$,
\deleted{\unsure{fermés}}\inserted{ouverts}.
% crops/150_scan013_P07.png: the following live fragment lies below syst. fond.; its attachment is uncertain.
\unsure{plus|plein} : $T$
(\unsure{ainsi que voisinages} tubulaires absolus de $Y$)}
% P07 scope check: the nearby overworked (6) and underlined main-text NB belong to P05/P06, not to this marginal note.

\textcircled{2} Soit $Y\lhook\joinrel\relbar X$ une inclusion fermée d'espaces
 topologiques. Un \uline{voisinage conique} de \uline{$Y$ dans} $X$
est un voisinage $T$ de $Y$ dans $X$
muni d'une opération (1) $I\times T\mathrel{\relbar\joinrel\relbar}T$
\inserted{faisant de $T$ un vois. \deleted{tub.}\inserted{coniq.}
absolu de $Y$.} \deleted{satisfaisant les conditions \unsure{dites.|id.}}
Alors les
$\lambda T=T_\lambda$ $(\lambda\in]0,1])$
forment d'autres voisinages
tubulaires de $Y$ dans $X$ (famille croissante) et
on voit que pour $\lambda\in]0,1[$, on a
\[
(\unsure{8})\qquad\dot T_\lambda=\lambda\dot T
\quad\text{est la frontière de }T_\lambda=\lambda T\text{ dans }X.
\]
\margin{12}
